% Imporditud: tools/import_eksperimendid.py
% Allikas: Eesti füüsikaolümpiaadi eksperimendi ülesannete
%   kogu (Rael Kalda, 2023), ülesanne Ü35, lahendus L35.
\setAuthor{EFO žürii}
\setRound{piirkonnavoor}
\setYear{1997}
\setNumber{G E2}
\setDifficulty{2}
\setTopic{termodünaamika}
\setVahendid{piirituslamp, kaalud, vihid, termomeeter, mõõtesilinder, alumiiniumist anum, vesi}
\setPunktid{12}
\setAllikas{Eksperimendi kogu Ü35, lahendus lk 48}
\setLahendusseis{toores}

\prob{Piirituse põlemine}
Kui suur osa piirituse põletamisel vabanevast soojushulgast kulub vee soojendamiseks? Hinnata mõõteviga. Piirituse kütteväärtus on 2,7 $\cdot$ 107J/kg, vee erisoojus 4200 J/kg $\cdot$ K, alumiiniumi erisoojus 880 J/kg $\cdot$ K.

\hint

\solu
Teooria: Soojendamise kasutegur: $\eta$ = (QA + QV )/Q , kus Q on piirituse põlemisel vabanenud soojushulk, QA anuma soojenemiseks kulunud soojushulk ja QV vee soojenemiseks kulunud soojushulk (2 p). Piirituse põlemisel vabanenud soojushulk: Q = rm, kus m on ära põlenud piirituse mass ja r on piirituse kütteväärtus (1 p). Vee soojenemiseks on kulunud soojushulk: QV = cm($t_2$ $-$ $t_1$) , kus m on soojendatava vee mass, c on vee erisoojus, $t_1$ vee algtemperatuur ja $t_2$ vee lõpptemperatuur (1 p). Samasuguse valemi järgi arvutatakse ka QA. Suhtelise vea valem: E(x) = $\Delta$x/x (1 p); Summa (vahe) arvutamise viga: E ($x_1$ $\pm$ $x_2$) = ($\Delta$$x_1$ + $\Delta$$x_2$) / ($x_1$ $\pm$ $x_2$) (1 p); Korrutise (jagatise) arvutamise viga: E ($x_1$/$x_2$) = E ($x_1$ $\cdot$ $x_2$) = E ($x_1$) + E ($x_2$) (1 p).

Mõõtmised: Mõõdame vee ja anuma massid (1 p). Mõõdame piirituslambi massi enne ja pärast põletamist (1 p). Mõõdame vee temperatuuri enne ja pärast soojendamist (1 p).

Mõõtevea arvutamine: Mõõtmiste suhtelised vead (1 p): Masside mõõtmiste suhteline viga E(m) = $\Delta$m/m, kus $\Delta$m on vastavalt vee või

anuma massi mõõtmise absoluutne viga (võib olla 0,1 g kuni 0,5 g sõltuvalt olemasolevast vihtide komplektist) ja m vastav mõõdetud massi väärtus. Temperatuuri vahede mõõtmiste suhteline viga E($\Delta$t) = 2$\Delta$t/($t_2$ $-$ $t_1$), kus $t_1$ ja $t_2$ on vastavalt mõõdetud vee ja anuma alg- ja lõpptemperatuuride väärtusud ning $\Delta$t on kasutatud termomeetri mõõtetäpsus (võib olla 0,1 $^\circ$ kuni 0,5 $^\circ$ sõltuvalt termomeetri skaalast).

Arvutuste suhtelised vead (1 p): Piirituse põlemisest eraldunud soojushulga suhteline viga on võrdne piirituslambi massi mõõtmise suhtelise veaga: E(Q) = E(m) = 2$\Delta$m/(m$^2$ $-$ $m_1$), kus $m_1$ ja m$^2$ on piirituslambi alg- ja lõppmassid ja $\Delta$m on massi mõõtmise absoluutne viga. Vee soojenemiseks kulunud soojushulga suhteline viga on võrdne vee massi ja temperatuuride vahe mõõtmiste suhteliste vigade summaga: E(QV ) = E(m) + E($\Delta$t). Anuma soojenemiseks kulunud soojushulga suhteline viga on võrdne anuma massi ja temperatuuride vahe mõõtmiste suhteliste vigade summaga: E(QA) = E(m)+ E($\Delta$t). Vee soojendamise kasuteguri suhteline viga on ülalarvutatud kolme soojushulga suhteliste vigade summa:E($\eta$) = E(Q) + E(QA) + E(QV ).
\probend
